% !TeX root = closed-systems-from-comparison-completeness.tex
% ============================================================
\chapter{Intrinsic Scalar Invariant and Discrete Hierarchy of Matter Sectors}
\label{ch:matter-spectral}
\label{ch:matter}
% ============================================================

% ============================================================
\section{Introduction and logical status}
\label{sec:spectral-intro}
% ============================================================

The theorem package
(\cref{thm:matter-spectrum-reduction,cor:matter-unique-forced-output})
establishes that every primitive excitation sector of the closed
comparison system \((U,\mathcal C)\) is classified, up to canonical
unitary equivalence, by irreducible representation data of the internal
gauge product
\[
  G_{\mathrm{int}} := U(1)\times SU(2)\times SU(3).
\]
The remaining classification problem is to distinguish sectors that
carry identical gauge representation data.
The present chapter constructs such a distinguishing invariant
intrinsically from the stabilized quadratic transport carrier
\[
  \mathcal K \cong F^2/F^3.
\]
Specifically, the chapter proves three things.
First, under the inherited second-jet faithfulness condition, the
restriction of the canonical scalar channel to the integral carrier is
identified, on the realized degree--\(2\) channel attached to the
stabilized quadratic carrier, with the scalar curvature projection:
\[
  Q|_{\mathcal K} = c \cdot \Scal(\cdot),\qquad c > 0.
\]
Second, the resulting scalar invariant $m_0(\sigma) := Q(\kappa_\sigma)$
is strictly positive and transport-invariant.
Third, the refinement depth $k_\sigma$ separates all primitive sectors
within a fixed gauge class, yielding a discrete hierarchy indexed by
$\mathbb N$.

The logical dependencies are as follows.
\begin{enumerate}[label=\textup{(\roman*)}]
\item
  Stabilization of the quadratic carrier and interface detection:
  \cref{thm:interface-stabilized,thm:interface-seal}.
\item
  Under the inherited second-jet faithfulness condition, nonzero
  stabilized square classes in the stabilized quadratic carrier are
  equivalent to nonzero realized curvature in any compatible smooth
  realization, and the scalar-channel argument below uses the realized
  degree--\(2\) channel attached to that carrier:
  \cref{prop:fusion-curvature,thm:interface-curvature}.
\item
  Isotropy completeness: \cref{thm:einstein-isotropy-complete}.
\item
  The admissible branch weight is additive on $\mathcal{K}$,
  $O(V,g_p)$-invariant, nonneg\-ative on admissible classes, and the
  admissible scalarization space is one-dimensional:
  \cref{def:admissible-branch-weight} condition~\textup{(S4)},
  \cref{lem:K-scalar-unique,lem:branch-weight-1d}.
\item
  The scalar summand of the curvature module is one-dimensional;
  the Weyl and trace-free Ricci summands carry no $O(V,g_p)$-invariant
  linear functional:
  \cref{sec:einstein-scalar,thm:einstein-beta-from-carrier}.
\item
  Canonical quadratic defect section, its basepoint-independence, and
  functorial uniqueness: \cref{thm:fusion-kappa}.
\item
  Existence, strict positivity, and quadratic extension of the scalar
  channel:
  \cref{thm:quadratic-scalar,lem:hilbert-strict-positivity,thm:quadratic-extension}.
\item
  Finite typing, compactness, and stability:
  \cref{thm:stability,thm:em-finite-phase-typing}.
\item
  Inverse-limit separation: \cref{thm:inverse-limit}.
\item
  Classification and finiteness of primitive transport mechanisms:
  \cref{thm:matter-primitive-uniqueness,thm:matter-no-proliferation}.
\item
  Primitive excitation dichotomy: \cref{thm:matter-excitation-dichotomy}.
\item
  No-third-locus exhaustion of representative enrichment:
  \cref{thm:representative-enrichment-classification}.
\end{enumerate}
No additional axiom, structural hypothesis, or dynamical assumption is
introduced beyond what is already established in Chapters~1--19.

% ============================================================
\section{Primitive mechanism package and spectrum reduction}
\label{sec:spectral-primitive-package}
% ============================================================

\begin{definition}[Primitive excitation mechanism]
\label{def:matter-primitive-mechanism}
A \emph{primitive excitation mechanism} is a nontrivial transport-visible
sector mechanism that is irreducible at its first finite visible stage,
in the sense that it is not an admissible composite of lower-stage
mechanisms of the same locus type.
\end{definition}

\begin{lemma}[Object-locus reduction]
\label{lem:matter-object-reduction}
Up to admissible equivalence, every primitive object-locus excitation
mechanism reduces to the canonical orientation-splitting doublet channel.
\end{lemma}

\begin{proof}
Let $\mathfrak m$ be a primitive object-locus excitation mechanism.
By the lift classification
(\cref{thm:lift-classification-main,cor:two-loci-lift}), object-locus
data are exactly representative-choice data.
At the primitive level, the only nontrivial object-locus first-visible
transport distinction is orientation splitting on the quadratic carrier.
Its existence and uniqueness up to admissible filtration-compatible
equivalence are given by
\cref{thm:orientation-splitting-uniqueness}.
Moreover, \cref{cor:orientation-primitive-doublet} identifies the
corresponding primitive irreducible block as a two-dimensional doublet.
Since $\mathfrak m$ is primitive
(\cref{def:matter-primitive-mechanism}), it must coincide with this
minimal object-locus primitive channel up to admissible equivalence.
Therefore every primitive object-locus mechanism reduces to the canonical
orientation-splitting doublet channel.
\end{proof}

\begin{lemma}[Morphism-locus reduction]
\label{lem:matter-morphism-reduction}
Up to admissible equivalence, every primitive morphism-locus excitation
mechanism reduces to minimal composite descent of index~$3$, i.e.
to the primitive triplet channel.
\end{lemma}

\begin{proof}
By
\cref{thm:em-minimal-composite-index,thm:em-denominator3-uniqueness},
the first nontrivial primitive morphism-locus descent has unique minimal
index~$3$.
Then \cref{cor:em-primitive-triplet} identifies the associated primitive
irreducible channel as a triplet, unique up to admissible equivalence.
\end{proof}

\begin{theorem}[Primitive mechanism uniqueness]
\label{thm:matter-primitive-uniqueness}
Primitive excitation mechanisms are exhausted, up to admissible
equivalence, by the following two primitive classes, possibly with
both and with no third primitive enrichment class:
\begin{enumerate}[label=\textup{(\roman*)}]
\item object-locus orientation splitting (doublet channel);
\item morphism-locus minimal composite descent (triplet channel).
\end{enumerate}
\end{theorem}

\begin{proof}
\textbf{Step~1 (exhaustion of possible loci).}
By the lift classification
(\cref{thm:lift-classification-main,cor:two-loci-lift}), primitive
mechanism data can occur only in representative-choice (object-locus)
or transport (morphism-locus) coordinates.
By
\cref{thm:representative-enrichment-classification}, there is no third
independent enrichment locus.
Hence every primitive excitation mechanism belongs to one of these two
locus types.

\textbf{Step~2 (uniqueness inside each locus).}
If a primitive mechanism is object-locus, then
\cref{lem:matter-object-reduction} reduces it, up to admissible
equivalence, to the canonical orientation-splitting doublet channel.
If a primitive mechanism is morphism-locus, then
\cref{lem:matter-morphism-reduction} reduces it, up to admissible
equivalence, to minimal composite descent of index~$3$, i.e. the
primitive triplet channel.
So each of the two loci contributes exactly one primitive type up to
admissible equivalence.

\textbf{Step~3 (exhaustion of primitive types).}
Step~1 shows that every primitive mechanism is exhausted by the object
and morphism locus forms with no third primitive enrichment class,
while Step~2 shows that each of those two classes has the listed
canonical primitive realization up to admissible equivalence.
Therefore primitive excitation mechanisms are exhausted by the two
listed classes: object-locus orientation splitting and morphism-locus
minimal composite descent.
\end{proof}

\begin{theorem}[No primitive proliferation]
\label{thm:matter-no-proliferation}
No primitive excitation mechanism exists beyond the two classes of
\cref{thm:matter-primitive-uniqueness}.
\end{theorem}

\begin{proof}
By \cref{thm:representative-enrichment-classification}, representative
enrichment is exhausted by object and morphism loci with no third source.
By \cref{thm:matter-primitive-uniqueness}, each locus has a unique
primitive mechanism type up to admissible equivalence.
Therefore no additional primitive mechanism can appear.
\end{proof}

\begin{theorem}[Primitive excitation two-class exhaustion]
\label{thm:matter-excitation-dichotomy}
Every nontrivial primitive excitation sector is exhausted by the two
primitive classes
\begin{enumerate}[label=\textup{(\roman*)}]
\item object-locus orientation-splitting class;
\item morphism-locus minimal composite-descent class.
\end{enumerate}
A sector may use either of these loci, or both, but no third
primitive enrichment class occurs.
\end{theorem}

\begin{proof}
By \cref{thm:lift-classification-main,cor:two-loci-lift}, every
transport-visible primitive sector is carried by object-locus data,
morphism-locus data, or both.
By \cref{thm:representative-enrichment-classification}, there is no
third independent enrichment locus.
By \cref{thm:matter-primitive-uniqueness}, the object locus
contributes only the canonical orientation-splitting doublet channel
and the morphism locus contributes only the primitive minimal
composite-descent triplet channel, up to admissible equivalence.
By \cref{thm:matter-no-proliferation}, no additional primitive
mechanism exists beyond those two locus-wise primitive classes.
Hence every nontrivial primitive excitation sector is exhausted by
the two listed primitive classes, with no third primitive
enrichment class.
\end{proof}

\begin{theorem}[Matter spectrum reduction]
\label{thm:matter-spectrum-reduction}
Under \cref{sp:closed-world}, primitive sector classification reduces to
irreducible representation data of
\[
  G_{\mathrm{int}} := U(1)\times SU(2)\times SU(3),
\]
with factors fixed intrinsically by:
\begin{enumerate}[label=\textup{(\roman*)}]
\item canonical phase rigidity for the abelian factor
      (\cref{thm:quadratic-phase-rigidity});
\item primitive object-locus doublet channel
      (\cref{cor:orientation-primitive-doublet});
\item primitive morphism-locus triplet channel
      (\cref{cor:em-primitive-triplet}).
\end{enumerate}
\end{theorem}

\begin{proof}
The phase factor is fixed by
\cref{thm:quadratic-phase-rigidity}, giving the canonical
one-dimensional abelian unitary channel.
By
\cref{lem:matter-object-reduction,lem:matter-morphism-reduction}, the
only primitive non-abelian multiplicity channels are the doublet and
triplet channels.
By \cref{thm:matter-no-proliferation}, no additional primitive channel is
available.
Therefore the primitive sector spectrum is exactly reduced to
representation data of the stated product.
\end{proof}

\begin{corollary}[Unique forced primitive output]
\label{cor:matter-unique-forced-output}
The primitive internal output is uniquely forced, up to admissible
equivalence, to the
\(1\)-, \(2\)-, and \(3\)-channel package encoded by
\(U(1)\times SU(2)\times SU(3)\).
\end{corollary}

\begin{proof}
By \cref{thm:matter-spectrum-reduction}, every primitive sector is
classified, up to admissible equivalence, by irreducible representation
data of
\[
  U(1)\times SU(2)\times SU(3),
\]
with the three factors fixed intrinsically by phase rigidity, primitive
object-locus splitting, and primitive morphism-locus descent.
By \cref{thm:matter-no-proliferation}, no additional primitive mechanism
exists beyond these forced channels.
Hence the primitive internal output is uniquely determined, up to
admissible equivalence, by the \(1\)-, \(2\)-, and \(3\)-channel
package.
\end{proof}

\begin{remark}[Role of the no-third-locus theorem]
\label{rem:spectral-two-locus-role}
The separation argument requires that every transport-distinguishing
datum between sectors factors through the two classified primitive
mechanisms.
By \cref{thm:representative-enrichment-classification}, every
representative enrichment of quotient semantics is exhausted by the
object locus and the morphism locus with no third independent source.
Under \cref{thm:matter-primitive-uniqueness,thm:matter-no-proliferation},
every such mechanism is enumerated and fully encoded at its minimal stage.
Hence any inter-sector distinction must be of a type already classified
and visible at the minimal stage.
\end{remark}

% ============================================================
\section{Orientation involution on the stabilized carrier}
\label{sec:spectral-orientation}
% ============================================================

The stabilized quadratic carrier inherits a canonical involution from
loop reversal, established in
\cref{sec:em-orientation-odd,sec:orientation-carrier}.

\begin{definition}[Orientation involution]
\label{def:spectral-orientation-involution}
The loop-reversal map $\tau(\gamma):=\gamma^{-1}$ on transport loops
induces, by the loop-reversal construction of
\cref{sec:em-orientation-odd},
a well-defined linear involution
\[
  \tau_{\mathcal K} : \mathcal K_{\mathbb R} \to \mathcal K_{\mathbb R},
  \qquad
  \tau_{\mathcal K}^2 = \mathrm{id},
\]
on $\mathcal K_{\mathbb R} := \mathcal K \otimes_{\mathbb Z} \mathbb R$.
\end{definition}

\begin{proposition}[Canonical eigenspace decomposition]
\label{prop:spectral-orientation-decomp}
There is a canonical orthogonal direct-sum decomposition
\[
  \mathcal K_{\mathbb R} = \mathcal K^+ \oplus \mathcal K^-,
  \qquad
  \mathcal K^\pm := \ker(\tau_{\mathcal K} \mp \mathrm{id}),
\]
into the $\pm 1$ eigenspaces of $\tau_{\mathcal K}$.
This coincides with the decomposition $H = H^+ \oplus H^-$ of
\cref{sec:em-orientation-odd}.
\end{proposition}

\begin{proof}
Since $\tau_{\mathcal K}^2 = \mathrm{id}$, the operator $\tau_{\mathcal K}$
is an orthogonal involution on $(\mathcal K_{\mathbb R}, g_{\mathcal K})$
(\cref{prop:hilbert-real-inner}) with minimal polynomial dividing
$(x-1)(x+1)$.
The standard spectral decomposition yields the stated splitting;
orthogonality follows because $\tau_{\mathcal K}$ is orthogonal.
The identification with $H^\pm$ is \cref{sec:em-orientation-odd}.
\end{proof}

\begin{definition}[Orientation parity]
\label{def:spectral-orientation-parity}
For a nonzero class $x \in \mathcal K_{\mathbb R}$, the
\emph{orientation parity} is
\[
  \epsilon(x) :=
  \begin{cases}
    +1 & x \in \mathcal K^+, \\
    -1 & x \in \mathcal K^-.
  \end{cases}
\]
\end{definition}

\begin{remark}[Parity is a $\mathbb Z_2$-grading]
\label{rem:parity-not-hierarchy}
Orientation parity records the bosonic/fermionic split of
\cref{thm:matter-excitation-dichotomy}; it does not separate sectors
within a single parity class.
The hierarchy invariant below is strictly finer.
\end{remark}

% ============================================================
\section{Defect class of an excitation sector}
\label{sec:spectral-defect}
% ============================================================

\begin{definition}[Defect class]
\label{def:spectral-defect-class}
Let $\sigma$ be a primitive excitation sector
(\cref{def:matter-primitive-mechanism}).
Fix a representative state $u_\sigma \in U$.
The \emph{defect class} of $\sigma$ is
\[
  \kappa_\sigma := \kappa(u_\sigma) \in \mathcal K,
\]
where $\kappa : U \to \mathcal K$ is the canonical quadratic defect
section of \cref{thm:fusion-kappa}.
\end{definition}

\begin{lemma}[Well-definedness]
\label{lem:spectral-defect-welldefined}
The class $\kappa_\sigma$ is independent of the choice of representative.
\end{lemma}

\begin{proof}
If $u_\sigma, u'_\sigma$ both represent $\sigma$, there exists
$g \in \Aut(U,\mathcal C)$ with $u'_\sigma = g \cdot u_\sigma$.
By equivariance of $\kappa$ (\cref{thm:fusion-kappa}),
$\kappa(u'_\sigma) = g_* \kappa(u_\sigma)$.
By the functorial uniqueness clause of \cref{thm:fusion-kappa},
$g_*$ preserves the class in $\mathcal K$, so
$\kappa(u'_\sigma) = \kappa(u_\sigma)$.
\end{proof}

\begin{lemma}[Nontriviality implies nonzero defect class]
\label{lem:spectral-defect-nonzero}
If $\sigma$ is a nontrivial primitive sector, then $\kappa_\sigma \neq 0$.
\end{lemma}

\begin{proof}
A nontrivial sector produces a nontrivial transport perturbation.
By \cref{thm:interface-stabilized,thm:interface-seal}, every nontrivial
transport obstruction has a nonzero representative in $\mathcal K$.
Hence $\kappa_\sigma \neq 0$.
\end{proof}

% ============================================================
\section{Identification of the scalar channel}
\label{sec:spectral-scalar-id}
% ============================================================

Under the inherited second-jet faithfulness condition, this section
identifies the restriction $Q|_{\mathcal K}$ on the realized
degree--\(2\) channel attached to the stabilized quadratic carrier
with the scalar curvature projection of the defect class.

\begin{remark}[Two levels of the scalar channel]
\label{rem:two-levels-Q}
The canonical scalar channel operates at two levels.
At the \emph{integral level}, $Q|_{\mathcal K} : \mathcal K \to
\mathbb R_{\geq 0}$ is a nonneg\-ative additive map on the abelian group
$\mathcal K$; nonneg\-ativity holds on admissible defect classes
(\cref{def:admissible-branch-weight}).
At the \emph{real level}, $Q_{\mathbb R} : \mathcal K_{\mathbb R} \to
\mathbb R_{\geq 0}$ is the unique positive quadratic extension
(\cref{thm:quadratic-scalar}).
\Cref{thm:spectral-scalar-identification} below concerns the integral
level; the form of $Q_{\mathbb R}$ is an open problem
(\cref{rem:spectral-scalar-id-scope}).
\end{remark}

\begin{remark}[Nonnegativity of scalar curvature on admissible classes]
\label{rem:scal-sign}
The scalar curvature $\Scal : \mathcal R(V) \to \mathbb R$ is
sign-indefinite on the full curvature module.
Under the inherited second-jet faithfulness condition, however,
admissible defect classes $\kappa\in\mathcal K$ are read on the
realized degree--\(2\) channel attached to the stabilized quadratic
carrier, and on that channel
\cref{thm:spectral-scalar-identification} gives
\[
  Q(\kappa) = c \cdot \Scal(\kappa),
  \qquad c>0.
\]
Since the admissible branch weight is nonnegative on admissible defect
classes (\cref{def:admissible-branch-weight}), it follows under that
inherited hypothesis that
\[
  \Scal(\kappa) \ge 0
  \qquad \text{for admissible } \kappa \in \mathcal K.
\]
This is not an additional axiom: it records the sign consequence of
the repaired scalar-channel identification on admissible classes, not
a positivity claim on the full curvature module.
\end{remark}

\begin{theorem}[Scalar channel identification]
\label{thm:spectral-scalar-identification}
Under the inherited second-jet faithfulness condition, the canonical
scalar channel
$Q : \mathcal K \to \mathbb R_{\geq 0}$ satisfies
\[
  Q(\kappa) = c \cdot \Scal(\kappa)
  \qquad (\kappa \in \mathcal K),
\]
for a universal constant $c > 0$, where
\[
  \Scal(\kappa) := \Phi \circ \pr_{\mathrm{scal}} \circ \Jtwo(\kappa),
\]
$\Jtwo : \mathcal K \to \mathcal R(V)$ is used only through
the realized degree--\(2\) channel attached to the stabilized
quadratic carrier under that inherited faithfulness condition,
$\pr_{\mathrm{scal}} : \mathcal R(V) \to \mathcal S(V)$ is the
$O(V,g_p)$-equivariant projection onto the one-dimensional scalar
summand (\cref{sec:einstein-scalar}), and
$\Phi : \mathcal S(V) \to \mathbb R$ is a fixed nonzero $O(V,g_p)$-invariant
linear functional.
\end{theorem}

\begin{proof}
\textbf{Step~1: $Q|_{\mathcal K}$ is additive and $O(V,g_p)$-invariant.}
By \cref{lem:branch-weight-1d}, $Q$ is a positive multiple of every
admissible branch weight.
By condition~\textup{(S4)} (\cref{def:admissible-branch-weight}), every
admissible branch weight is additive on $\mathcal K$ as a group
homomorphism and natural under the realized symmetry flow.
Hence $Q|_{\mathcal K}$ is additive on $\mathcal K$ and invariant under
the realized local symmetry group.
By \cref{thm:einstein-isotropy-complete}, the realized local symmetry
group at each point $p$ is $O(V,g_p)$.
Therefore, under the inherited second-jet faithfulness condition, the
transported scalar channel on the realized degree--\(2\) channel
attached to the stabilized quadratic carrier is an
$O(V,g_p)$-invariant nonnegative additive map.

\textbf{Step~2: Passage from the stabilized quadratic carrier to its realized degree--\(2\) channel.}
By \cref{prop:fusion-curvature,thm:interface-curvature}, the map
$\Jtwo : \mathcal K \to \mathcal R(V)$ carries the stabilized
quadratic carrier into the realized degree--\(2\) channel attached to
that carrier, where nonzero stabilized square classes are equivalent
to nonzero realized curvature.

\textbf{Step~3: $O(V,g_p)$-invariant additive maps on $\mathcal R(V)$
span a one-dimensional space.}
The $O(V,g_p)$-irreducible decomposition
\[
  \mathcal R(V) = \mathcal W(V) \oplus \mathcal E(V) \oplus \mathcal S(V)
\]
is established in \cref{sec:einstein-scalar};
$\mathcal S(V)$ is one-dimensional (\cref{thm:einstein-beta-from-carrier}).
Since $\mathcal W(V)$ and $\mathcal E(V)$ are non-trivial irreducible
$O(V,g_p)$-representations,
\[
  \operatorname{Hom}_{O(V,g_p)}\!\bigl(\mathcal W(V),\,\mathbb R\bigr) = 0,
  \qquad
  \operatorname{Hom}_{O(V,g_p)}\!\bigl(\mathcal E(V),\,\mathbb R\bigr) = 0.
\]
Every $O(V,g_p)$-invariant additive map $\mathcal R(V) \to \mathbb R$
therefore annihilates $\mathcal W(V) \oplus \mathcal E(V)$ and factors
through $\pr_{\mathrm{scal}}$.
The space of such maps is one-dimensional, spanned by
$\Phi \circ \pr_{\mathrm{scal}}$.

\textbf{Step~4: Conclusion.}
By Steps~1--3, that transported scalar channel on the realized
degree--\(2\) channel is
$c' \cdot (\Phi \circ \pr_{\mathrm{scal}})$
for some $c' \in \mathbb R$.
Since $Q$ is nonnegative and nontrivial on admissible alternatives
(\cref{lem:hilbert-strict-positivity}), we have $c' > 0$.
Set $c := c'$; pulling this identity back along that realized
degree--\(2\) channel gives the stated formula.
\end{proof}

\begin{remark}[Scope of the identification]
\label{rem:spectral-scalar-id-scope}
\Cref{thm:spectral-scalar-identification} identifies $Q|_{\mathcal K}$
as an additive map on the integral carrier $\mathcal K$.
The symmetric bilinear source of the quadratic extension is forced by
polarization (\cref{thm:scalar-pairing}).  What remains open here is not
the existence of a bilinear source, but the fully explicit real-level formula
relating
$Q_{\mathbb R} : \mathcal K_{\mathbb R} \to \mathbb R_{\geq 0}$
to $\Scal$ beyond the integral carrier.  Determining that explicit formula
for $Q_{\mathbb R}$ remains an open problem.
\end{remark}

\begin{remark}[Relation to the cosmological scalar]
\label{rem:spectral-cosmological}
Under the inherited second-jet faithfulness condition, the map
$Q|_{\mathcal K} = c \cdot \Scal(\cdot)$
is the same scalar projection as the cosmological coefficient
$\beta_\triangle$ in \cref{thm:einstein-beta-from-carrier}.
Consequently, under that inherited hypothesis, the Born probability
law (\cref{thm:born}) and the cosmological scalar term arise from the
same scalar projection on the realized degree--\(2\) channel attached
to the stabilized quadratic carrier.
\end{remark}

% ============================================================
\section{The canonical scalar invariant}
\label{sec:spectral-scalar}
% ============================================================

\begin{definition}[Scalar invariant]
\label{def:spectral-scalar}
For a primitive excitation sector $\sigma$, the \emph{scalar invariant} is
\[
  m_0(\sigma) := Q(\kappa_\sigma),
\]
where $Q : \mathcal K_{\mathbb R} \to \mathbb R_{\geq 0}$ is the
canonical quadratic scalar channel of \cref{thm:quadratic-scalar}.
Under the inherited second-jet faithfulness condition, on the integral
carrier
$m_0(\sigma) = c \cdot \Scal(\kappa_\sigma)$
by \cref{thm:spectral-scalar-identification}.
\end{definition}

\begin{theorem}[Strict positivity]
\label{thm:spectral-scalar-positive}
For every nontrivial primitive sector $\sigma$, $m_0(\sigma) > 0$.
\end{theorem}

\begin{proof}
By \cref{lem:spectral-defect-nonzero}, $\kappa_\sigma \neq 0$.
By \cref{lem:hilbert-strict-positivity}, $Q(x) > 0$ for every nonzero
$x \in \mathcal K_{\mathbb R}$.
Hence $m_0(\sigma) = Q(\kappa_\sigma) > 0$.
\end{proof}

\begin{theorem}[Transport invariance]
\label{thm:spectral-scalar-invariant}
The scalar $m_0(\sigma)$ is invariant under admissible transport.
\end{theorem}

\begin{proof}
Let $\Phi_t \in \Aut(U,\mathcal C)$ and $u \in U$.
By equivariance of the canonical quadratic defect section
(\cref{thm:fusion-kappa}),
\[
  \kappa(\Phi_t(u)) = \Phi_{t*}\kappa(u).
\]
By \cref{thm:quadratic-scalar}, the restriction of $Q$ to the integral
carrier is the canonical degree--\(2\) scalar channel, and by the
naturality clause in condition~\textup{(S4)} of
\cref{def:admissible-branch-weight} that scalar channel is preserved by
admissible symmetry flow on the realized degree--\(2\) channel attached
to the stabilized quadratic carrier. Therefore
\[
  Q(\Phi_{t*}k)=Q(k)
  \qquad
  (k\in \mathcal K).
\]
Applying this to $k=\kappa(u)$ gives
\[
  Q(\kappa(\Phi_t(u)))
  =
  Q(\Phi_{t*}\kappa(u))
  =
  Q(\kappa(u)).
\]
Hence $m_0(\sigma) = Q(\kappa_\sigma)$ is preserved under admissible
transport.
\end{proof}

\begin{remark}[Non-injectivity of $m_0$]
\label{rem:m0-not-injective}
Under the inherited second-jet faithfulness condition, if two distinct
sectors have the same scalar-curvature coordinate of their defect
class, then they share the same value of $m_0$.
The separation theorem uses the strictly finer invariant $k_\sigma$.
\end{remark}

% ============================================================
\section{Refinement depth}
\label{sec:spectral-depth}
% ============================================================

The stabilized carrier is $\mathcal K = \varprojlim_k \mathcal K_k$
(\cref{thm:interface-stabilized}), with projections
$\pi_k : \mathcal K \to \mathcal K_k$.

\begin{definition}[Refinement depth]
\label{def:spectral-depth}
For a nontrivial primitive sector $\sigma$, the
\emph{refinement depth} is
\[
  k_\sigma
  := \min\bigl\{ k \in \mathbb N : \pi_k(\kappa_\sigma) \neq 0 \bigr\}.
\]
\end{definition}

\begin{proposition}[Finiteness of refinement depth]
\label{prop:spectral-depth-finite}
Every nontrivial primitive sector has $k_\sigma < \infty$.
\end{proposition}

\begin{proof}
By \cref{lem:spectral-defect-nonzero}, $\kappa_\sigma \neq 0$ in
$\varprojlim_k \mathcal K_k$.
By \cref{thm:inverse-limit}, at least one stage projection is nonzero,
so the minimum exists in $\mathbb N$.
\end{proof}

\begin{proposition}[Transport invariance]
\label{prop:spectral-depth-invariant}
The refinement depth $k_\sigma$ is invariant under admissible transport.
\end{proposition}

\begin{proof}
Automorphisms of $(U,\mathcal C)$ preserve the projective system
$\{\mathcal K_k\}$ and commute with the projections $\pi_k$ by
functoriality of the defect construction (\cref{thm:fusion-kappa}).
Hence the minimal nonzero stage index is preserved.
\end{proof}

% ============================================================
\section{Finite-stage determination and stage rigidity}
\label{sec:spectral-rigidity}
% ============================================================

\begin{lemma}[Finite-stage separation]
\label{lem:spectral-finite-sep}
Let $\sigma_1, \sigma_2$ be distinct primitive sectors.
Then there exists $k \in \mathbb N$ with
$\pi_k(\kappa_{\sigma_1}) \neq \pi_k(\kappa_{\sigma_2})$.
\end{lemma}

\begin{proof}
By \cref{thm:inverse-limit}, elements of $\varprojlim_k \mathcal K_k$
are equal if and only if all stage projections agree.
Contrapositive: $\kappa_{\sigma_1} \neq \kappa_{\sigma_2}$ implies the
existence of such $k$.

We show $\sigma_1 \neq \sigma_2$ forces $\kappa_{\sigma_1} \neq
\kappa_{\sigma_2}$.
By \cref{thm:lift-classification-main}, all transport distinction is
carried by the morphism locus.
By \cref{thm:em-finite-phase-typing}, any such distinction is typed by
finite comparison data in some stage-$k$ Boolean algebra $\mathcal B_k$.
If two sectors agreed at every finite stage they would agree in the
inverse limit by \cref{thm:stability,thm:em-closedness-phase},
contradicting inequivalence.
\end{proof}

\begin{lemma}[Stage rigidity]
\label{lem:spectral-stage-rigidity}
Let $\sigma_1, \sigma_2$ be distinct primitive sectors with identical
gauge data and $k_{\sigma_1} = k_{\sigma_2} = k$.
If $\pi_k(\kappa_{\sigma_1}) = \pi_k(\kappa_{\sigma_2})$, then
$\sigma_1 = \sigma_2$.
\end{lemma}

\begin{proof}
Assume $\pi_k(\kappa_{\sigma_1}) = \pi_k(\kappa_{\sigma_2})$.
Both defect classes vanish below stage $k$ by minimality and coincide at
stage $k$ by assumption.
Suppose for contradiction that $\sigma_1 \neq \sigma_2$.
By \cref{lem:spectral-finite-sep}, there exists $k' > k$ at which the
projections differ; new distinguishing data must appear above stage $k$.

\textbf{Step~1.}
By \cref{thm:representative-enrichment-classification}, every
representative enrichment is exhausted by the object locus and the
morphism locus.
Any inter-sector distinction must lie in one of these two channels.

\textbf{Step~2.}
By \cref{thm:matter-primitive-uniqueness}, the primitive mechanisms are
exactly object-locus orientation splitting
(\cref{lem:matter-object-reduction}) and morphism-locus minimal
composite descent of index~$3$ (\cref{lem:matter-morphism-reduction}).
By \cref{thm:matter-no-proliferation}, no further primitive mechanism
exists.
Each is characterized up to admissible equivalence by
\cref{thm:orientation-splitting-uniqueness,thm:em-denominator3-uniqueness}
and is fully encoded at its minimal stage: orientation splitting by the
stage-$k$ class of \cref{cor:orientation-primitive-doublet};
index-$3$ descent by the stage-$k$ class of \cref{cor:em-primitive-triplet}.

\textbf{Step~3.}
A new distinguishing datum at stage $k' > k$ would require either a new
primitive mechanism (excluded by \cref{thm:matter-no-proliferation}) or a
new enrichment locus (excluded by
\cref{thm:representative-enrichment-classification}).
Hence all finite stages agree, so $\sigma_1 = \sigma_2$ by
\cref{thm:stability,thm:em-closedness-phase}, a contradiction.
\end{proof}

% ============================================================
\section{Separation of sectors with identical gauge labels}
\label{sec:spectral-separation}
% ============================================================

\begin{theorem}[Separation by refinement depth]
\label{thm:spectral-depth-separation}
Let $\sigma_1, \sigma_2$ be nontrivial primitive sectors carrying
identical gauge representation data for
$U(1)\times SU(2)\times SU(3)$.
If $\sigma_1 \neq \sigma_2$, then $k_{\sigma_1} \neq k_{\sigma_2}$.
\end{theorem}

\begin{proof}
Suppose for contradiction that $k_{\sigma_1} = k_{\sigma_2} = k$.

\textbf{Case~1:} $\pi_k(\kappa_{\sigma_1}) \neq \pi_k(\kappa_{\sigma_2})$.
By \cref{thm:matter-excitation-dichotomy}, the locus type of a sector is
determined by its gauge labels.
By \cref{thm:matter-primitive-uniqueness}, at most one primitive mechanism
exists per locus type within a fixed gauge class.
Two distinct sectors sharing gauge labels, locus type, and depth $k$
would require two primitive mechanisms of the same locus type in the same
gauge class, contradicting uniqueness.

\textbf{Case~2:} $\pi_k(\kappa_{\sigma_1}) = \pi_k(\kappa_{\sigma_2})$.
By \cref{lem:spectral-stage-rigidity}, $\sigma_1 = \sigma_2$,
contradicting the assumption.

Both cases yield contradictions.
\end{proof}

\begin{corollary}[Discrete hierarchy within each gauge class]
\label{cor:spectral-hierarchy}
Within any fixed gauge representation class, the assignment
$\sigma \mapsto k_\sigma$ is injective on primitive sectors;
in particular, they are indexed by a discrete subset of $\mathbb N$.
\end{corollary}

\begin{proof}
Injectivity is \cref{thm:spectral-depth-separation}.
Discreteness holds since $k_\sigma \in \mathbb N$.
\end{proof}

% ============================================================
\section{Quadratic channel count at each depth}
\label{sec:spectral-dimension}
% ============================================================

\begin{theorem}[Quadratic channel count {\normalfont(\cref{thm:binomk2})}]
\label{thm:spectral-binomial}
For a free $k$-generator triangle-closed system,
\[
  \mathrm{rank}(\mathcal K_k) = \binom{k}{2},
  \qquad
  \mathcal K_k \cong \Lambda^2(\mathbb Z^k).
\]
\end{theorem}

\begin{remark}[Dimension versus value]
\label{rem:spectral-dimension-not-value}
\Cref{thm:spectral-binomial} gives the rank of $\mathcal K_k$.
By \cref{thm:spectral-scalar-identification}, under the inherited
second-jet faithfulness condition the scalar channel on admissible
defect classes carried by \(\mathcal K\) agrees, through the
realized degree--\(2\) channel attached to the stabilized quadratic
carrier, with the scalar curvature projection there.
The channel count constrains how many independent defect classes exist at
each depth; that conditional scalar-channel identification
constrains their values.
\end{remark}

% ============================================================
\section{Normalized ratio invariant}
\label{sec:spectral-ratio}
% ============================================================

The scalar channel $Q$ is unique up to positive normalization
(\cref{thm:quadratic-scalar}), so absolute values of $m_0(\sigma)$
depend on a normalization choice.
The normalization-independent datum is the ratio.

\begin{definition}[Ratio invariant]
\label{def:spectral-ratio}
Fix a nontrivial reference sector $\sigma_*$.
The \emph{ratio invariant} is
\[
  \widehat{m}(\sigma;\sigma_*)
  := \frac{m_0(\sigma)}{m_0(\sigma_*)}
  = \frac{Q(\kappa_\sigma)}{Q(\kappa_{\sigma_*})}.
\]
\end{definition}

\begin{theorem}[Properties of the ratio invariant]
\label{thm:spectral-ratio-properties}
The ratio $\widehat{m}(\sigma;\sigma_*)$ is
\begin{enumerate}[label=\textup{(\roman*)}]
\item independent of the normalization of $Q$;
\item invariant under admissible transport; and
\item strictly positive for every nontrivial $\sigma$.
\end{enumerate}
\end{theorem}

\begin{proof}
\textup{(i)}: If $Q' = \lambda Q$ for $\lambda > 0$, numerator and
denominator scale by $\lambda$, leaving the ratio unchanged.
\textup{(ii)}: Both $Q(\kappa_\sigma)$ and $Q(\kappa_{\sigma_*})$ are
transport-invariant by \cref{thm:spectral-scalar-invariant}.
\textup{(iii)}: The numerator satisfies $m_0(\sigma) > 0$ by
\cref{thm:spectral-scalar-positive}; the denominator satisfies
$m_0(\sigma_*) > 0$ by the same theorem.
\end{proof}

\begin{remark}[Relation to the appendix normalized scalar]
\label{rem:spectral-ratio-appendix}
The ratio $\widehat{m}(\sigma;\sigma_*)$ is the sector-level analogue of
the appendix normalized scalar
$\widehat{\mathfrak s}_{\mathrm{fus}}(u;u_*) = Q(\kappa(u))/Q(\kappa(u_*))$ of
\cref{sec:fusion-normalized}: both cancel the normalization freedom in
$Q$ and are preserved by admissible automorphisms.
\end{remark}

% ============================================================
\section{Conclusion}
\label{sec:spectral-conclusion}
% ============================================================

The canonical invariant that distinguishes primitive excitation sectors
with identical $U(1)\times SU(2)\times SU(3)$ gauge labels is the pair
\[
  \sigma \longmapsto
  \bigl(\epsilon(\kappa_\sigma),\, k_\sigma\bigr)
  \in \{+1,-1\} \times \mathbb N,
\]
where $\epsilon(\kappa_\sigma)$ is the orientation parity
(\cref{def:spectral-orientation-parity}) and $k_\sigma$ is the
refinement depth (\cref{def:spectral-depth}).

Orientation parity reproduces the bosonic/fermionic split of
\cref{thm:matter-excitation-dichotomy}.
Refinement depth separates distinct sectors within a fixed gauge class
by \cref{thm:spectral-depth-separation,cor:spectral-hierarchy}.

The scalar invariant $m_0(\sigma) = Q(\kappa_\sigma)$ is strictly
positive (\cref{thm:spectral-scalar-positive}) and transport-invariant
(\cref{thm:spectral-scalar-invariant}).
Under the inherited second-jet faithfulness condition, its restriction
to the integral carrier satisfies
$m_0(\sigma) = c \cdot \Scal(\kappa_\sigma)$
by \cref{thm:spectral-scalar-identification}.
The normalization-free version is $\widehat{m}(\sigma;\sigma_*)$
(\cref{def:spectral-ratio,thm:spectral-ratio-properties}).

The quadratic channel count $\mathrm{rank}(\mathcal K_k) = \binom{k}{2}$
(\cref{thm:spectral-binomial}) gives the dimension of the space of
admissible defect classes at each depth.

\[
\boxed{
\begin{aligned}
&\text{Gauge labels determine sector type.} \\[2pt]
&\text{Refinement depth } k_\sigma \text{ separates distinct sectors}\\
&\quad\text{within a fixed gauge class.} \\[2pt]
&\text{Under the inherited second-jet faithfulness condition,} \\
&Q|_{\mathcal K} = c \cdot \Scal(\cdot),\quad c > 0.
\end{aligned}
}
\]